\documentclass[11pt,a4paper]{article}
\usepackage[T1]{fontenc}
\usepackage[utf8]{inputenc}
\usepackage[main=italian,provide=*]{babel}
\usepackage{amsmath,amssymb}
\usepackage{geometry}
\geometry{margin=2.3cm,top=2.5cm,bottom=2.7cm}
\usepackage{booktabs}
\usepackage[table,dvipsnames]{xcolor}
\usepackage{tcolorbox}
\tcbuselibrary{skins,breakable}
\usepackage{titlesec}
\usepackage{enumitem}
\usepackage{needspace}
\usepackage{fancyhdr}
\usepackage{graphicx}
\usepackage{hyperref}
\hypersetup{colorlinks=true,linkcolor=Sepia,citecolor=Sepia,urlcolor=Sepia}
\newcommand{\Tr}{\operatorname{Tr}}
\newcommand{\ket}[1]{\lvert #1 \rangle}
\newcommand{\bra}[1]{\langle #1 \rvert}

\definecolor{inkblue}{HTML}{1B3A5C}
\definecolor{lightblue}{HTML}{EAF1F8}
\definecolor{accent}{HTML}{B0562A}
\definecolor{softgray}{HTML}{6B6B6B}
\definecolor{okgreen}{HTML}{2E6B3E}
\definecolor{okgreenbg}{HTML}{EAF5EC}

\titleformat{\section}
  {\normalfont\Large\bfseries\color{inkblue}}
  {\thesection}{0.6em}{}[{\color{inkblue!40}\titlerule}]
\titlespacing{\section}{0pt}{0.9em}{0.4em}
\titleformat{\subsection}
  {\normalfont\large\bfseries\color{accent}}
  {}{0em}{}
\titlespacing{\subsection}{0pt}{0.6em}{0.2em}

\pagestyle{fancy}
\fancyhf{}
\renewcommand{\headrulewidth}{0.4pt}
\fancyhead[L]{\small\color{softgray}Correlazioni dinamiche sotto rumore (anello)}
\fancyhead[R]{\small\color{softgray}\thepage}
\fancyfoot[C]{\small\color{softgray}Tesi triennale, Samuele --- Relatore: Prof.\ Alessandro Chiesa}

\newtcolorbox{keyeq}[1][]{
  colback=lightblue, colframe=inkblue!70, boxrule=0.6pt,
  arc=2pt, left=8pt, right=8pt, top=3pt, bottom=3pt, #1
}
\newtcolorbox{okbox}[1]{
  colback=okgreenbg, colframe=okgreen!60, boxrule=0.6pt,
  arc=2pt, left=7pt, right=7pt, top=2pt, bottom=2pt,
  fonttitle=\bfseries\color{okgreen}, title={#1}
}

\title{\vspace{-1.5em}\color{inkblue}Correlazioni dinamiche sotto rumore --- trimero ad anello\\[0.2em]
\Large\normalfont\color{softgray}Circuito Hadamard test a 4 qubit, rumore su tutti gli ingredienti}
\author{Note di lavoro --- Tesi triennale, Samuele\\ Relatore: Prof.\ Alessandro Chiesa}
\date{}

\begin{document}
\sloppy
\maketitle
\thispagestyle{fancy}
\vspace{-2.0em}

File di codice:
\texttt{correlatori\_\allowbreak rumorosi\_trimero\_anello.py} (mirror di
\texttt{correlatori\_rumorosi\_dimero.py}). Punto di lavoro: $J=1,\,J'=0.4,\,b=b_c=2.4,\,D=0.15$
--- stesso punto usato per l'esplorazione dei correlatori in Parte 1
(\texttt{trimero\_\allowbreak anello\_correlazioni.tex}), non il punto
$R_0$ (quello serve solo per lo Scenario B dello Stadio 3, mai usato per i
correlatori in Parte 1). Correlatore usato come riferimento in questo
documento: $C_{21}^{xx}(t)$, coerente con la scelta del dimero.

\tableofcontents

\section{Obiettivo}

Stesso circuito Hadamard test a 4 qubit di Parte 1 (3 di registro + 1
ancilla), ora con rumore su \textbf{tutti} i suoi ingredienti:
preparazione (Stadio 2), passo esterno di Trotter ripetuto $N$ volte
(Stadio 3, blocco riusato senza reimplementarlo), e infine la misura
sull'ancilla (errore di lettura, Stadio 1 --- mai agganciato finora
perché gli Stadi 2-3 non misuravano nulla).

\begin{keyeq}[title={Convenzione qubit: diversa dal dimero, ereditata da Parte 1}]
Registro = qubit 0,1,2; ancilla = qubit 3 --- \textbf{non} ancilla=0 come
sul dimero. Eredità diretta di
\texttt{circuito\_\allowbreak correlazioni\_\allowbreak trimero\_anello.py}
(Parte 1): non ridefinita qui, per evitare un disallineamento con il
circuito già validato. Conseguenza verificata esplicitamente (non
assunta): nel Monte Carlo di verifica (Sez.~\ref{sec:readout}), il bit
dell'ancilla nella stringa di misura è il carattere più a
\textbf{sinistra}, non più a destra come sul dimero.
\end{keyeq}

\section{Quattro blocchi, transpilazione a blocchi (Stadio 3 riusato)}

Stessa strategia already fissata allo Stadio 3, estesa qui a due nuovi
blocchi locali a 2 qubit:
\begin{enumerate}[leftmargin=1.4em,itemsep=0.2em]
\item preparazione (ansatz $W$-2q.6, Stadio 2) --- transpilata una volta;
\item controlled-$W$ sul sito $j$ (ancilla + 1 qubit di registro) ---
transpilata una volta;
\item passo esterno di Trotter (3 qubit), ripetuto $N$ volte ---
\textbf{riusato da} \texttt{trotter\_\allowbreak rumoroso\_trimero\_
anello.build\_step\_block} (Stadio 3), non reimplementato;
\item anti-controlled-$V$ + rotazione di base (ancilla + 1 qubit di
registro) --- transpilata una volta.
\end{enumerate}

\section{Risultati: limite di rumore nullo e conteggio gate}

\begin{center}
\renewcommand{\arraystretch}{1.25}
\begin{tabular}{@{}lccc@{}}
\toprule
\textbf{Caso $(i,\alpha,j,\beta,t,N)$} & \textbf{Parte 1 (ansatz)} & \textbf{Stadio 4 (rumore=0)} & $|$\textbf{diff}$|$\\
\midrule
$(2,x,1,x,2.0,20)$ & $-0.571048+0.361219i$ & $-0.571048+0.361219i$ & $3.5\times10^{-9}$\\
$(1,z,1,z,1.0,20)$ & $+0.398173+0.185174i$ & $+0.398173+0.185174i$ & $2.1\times10^{-9}$\\
$(3,x,3,y,0.5,20)$ & $-0.770489+0.308875i$ & $-0.770489+0.308875i$ & $2.1\times10^{-11}$\\
\bottomrule
\end{tabular}
\end{center}

\begin{okbox}{Verificato}
Coincidenza con Parte 1 entro l'infedeltà nota della preparazione VQE
(scarti $\sim10^{-9}$--$10^{-11}$), su tre combinazioni diverse di siti e
componenti (incluso un caso a stesso sito, $i=j=3$, componenti miste
$x,y$ --- non strutturalmente nullo per simmetria, a differenza delle
combinazioni $C_{33}^{xz},C_{33}^{zx},C_{33}^{yz},C_{33}^{zy}$ già
identificate come nulle in Parte 1).
\end{okbox}

Conteggio gate: incremento costante di \textbf{15 CNOT} per unità di $N$
(lo stesso costo per passo già stabilito allo Stadio 3), confermato da
$N=1$ a $N=40$ --- nessuna fusione dei passi.

\section{Formula del readout: verifica Monte Carlo indipendente}
\label{sec:readout}

Stessa formula analitica del dimero (verificata nello Stadio 1 come
indipendente dal numero di qubit): $\langle Z\rangle_\text{readout} =
(1-2p)\langle Z\rangle_\text{ideale}$, applicata dopo il rumore di gate.
Verificata qui con una simulazione Monte Carlo indipendente (\texttt{ReadoutError}
vero, $2\times10^5$ shot, $p=0.05$):

\begin{center}
\renewcommand{\arraystretch}{1.25}
\begin{tabular}{@{}lc@{}}
\toprule
\textbf{Quantità} & \textbf{Valore}\\
\midrule
$\langle Z\rangle$ analitico (gate noise + correzione readout) & $-0.232313$\\
$\langle Z\rangle$ Monte Carlo ($2\times10^5$ shot) & $-0.234790$\\
$|$differenza$|$ & $2.5\times10^{-3}$ (atteso $\lesssim6.7\times10^{-3}$, $3\sigma$)\\
\bottomrule
\end{tabular}
\end{center}

\begin{okbox}{Verificato}
Formula analitica confermata da un cammino di calcolo indipendente,
entro la tolleranza statistica attesa. Posizione del bit dell'ancilla
nella stringa di misura verificata esplicitamente (carattere più a
sinistra, qubit 3), non assunta dall'analogia con il dimero.
\end{okbox}

\section{Il compromesso Trotter/rumore sul correlatore}

\begin{figure}[h]
\centering
\includegraphics[width=0.82\linewidth]{figure/fig_correlatore_rumoroso_trimero_anello.pdf}
\caption{\textit{$|C_{21}^{xx}(t{=}2,N)|$ a rumore nullo e di riferimento
(\texttt{ibm\_torino}). A rumore nullo il correlatore converge rapidamente
(già stabile da $N\gtrsim10$); sotto rumore emerge un massimo a $N$
finito.}}
\end{figure}

\begin{okbox}{Verificato}
$N^*=3$, $|C(N^*)|=0.505$ --- diverso sia da $N^*_A=2$ (fedeltà di
Trotter, Stadio 3, stesso punto VQE) sia da $N^*_B=10$ (fedeltà, punto
$R_0$): tre osservabili diverse, tre $N^*$ diversi, coerente col
principio già stabilito sul dimero (``il numero ottimale di passi
dipende dall'osservabile che si sta calcolando'') e qui riverificato
direttamente su una quarta combinazione punto/osservabile, non assunto
per analogia.
\end{okbox}

\section{Cosa non è ancora stato controllato qui}

Nessuno scan sistematico sui parametri di rumore (Stadio 5). Il
correlatore è stato verificato per un solo punto di lavoro (quello VQE,
coerente con Parte 1); lo scenario $R_0$ non è stato esplorato per i
correlatori, né qui né in Parte 1 --- non è un'omissione di questo
documento, è coerente con lo scope già fissato quando Parte 1 ha
esplorato i correlatori solo al punto VQE.

\section{Considerazioni e cosa serve al passo successivo}

Lo Stadio 5 (scan sui parametri di rumore) riprenderà sia la fedeltà di
Trotter (Stadio 3, entrambi gli scenari) sia il correlatore (qui), per
vedere come i tre $N^*$ già trovati ($2$, $10$, $3$) si spostano al
variare di $(\varepsilon_{1q},\varepsilon_{2q},p_\text{readout})$. Sul
dimero, il modello di scomposizione additiva
$\varepsilon_\text{tot}(N)\sim a/N+g\varepsilon_{2q}N$ non descriveva
correttamente lo scaling di $N^*$ per una fedeltà (pendenza log-log
osservata $-0.195$ contro $-0.5$ prevista, dichiarato esplicitamente nel
documento dello Stadio 5 del dimero) --- non c'è quindi motivo di
aspettarsi che descriva meglio né la fedeltà né il correlatore
dell'anello: lo scan andrà letto sui dati, non sul modello.

\begin{okbox}{Sintesi in una riga}
Il circuito Hadamard test dell'anello, reso rumoroso riusando il blocco
di Trotter dello Stadio 3, riproduce Parte 1 nel limite di rumore nullo
ed espone un compromesso Trotter/rumore con $N^*=3$ --- un terzo valore
di $N^*$, diverso dai due già trovati allo Stadio 3.
\end{okbox}

\needspace{4\baselineskip}
\section*{File prodotti}
\addcontentsline{toc}{section}{File prodotti}

\texttt{correlatori\_\allowbreak rumorosi\_trimero\_anello.py},
\texttt{validate\_\allowbreak correlatori\_\allowbreak rumorosi\_trimero\_anello.py},
\texttt{generate\_figure\_\allowbreak correlatore\_\allowbreak rumoroso\_trimero\_anello.py},
questo documento.

\end{document}
